\section{Probabilities needed for disjunction}\label{appendix:calculations-prob-or}

Let $\{\TF^1_n\}_{n \in \mathbb{N}}$ be the family of probabilistic transformers recognizing $\gL^1$ and $\{\TF^2_n\}_{n \in \mathbb{N}}$ the family recognizing $\gL^2$.
Let $p_1$ be the probability of $\{\TF^1_n\}_{n \in \mathbb{N}}$ answering correctly and $p_2$ the probability of $\{\TF^2_n\}_{n \in \mathbb{N}}$ answering correctly.

We construct the family $\{\TF_n\}_{n \in \mathbb{N}}$ that runs $\{\TF^1_n\}_{n \in \mathbb{N}}$ and $\{\TF^2_n\}_{n \in \mathbb{N}}$ and outputs its disjunction.

What is the probability of $\{\TF_n\}_{n \in \mathbb{N}}$ answering correctly $\gL^1(\alpha) \lor \gL^2(\alpha)$?
For each word $\alpha$, there are four possible cases
($\alpha \in    \gL^1 \land \alpha \in    \gL^2,
  \alpha \in    \gL^1 \land \alpha \notin \gL^2,
  \alpha \notin \gL^1 \land \alpha \in    \gL^2,
  \alpha \notin \gL^1 \land \alpha \notin \gL^2$).
In Figure \ref{fig:cases}, we show the probability of $\TF^1$ and $\TF^2$ accepting or rejecting $\alpha$ in each of the cases.

\begin{figure}[htbp]
  \centering
  % Fila 1
  \begin{subfigure}[t]{0.48\linewidth}
    \centering
    \resizebox{\linewidth}{!}{%
    \begin{tabular}{|r|c|c|}
      \hline
                      & $\TF^1$ accepts   & $\TF^1$ rejects \\
      \hline
      $\TF^2$ accepts & $p_1p_2$          & $(1-p_1)p_2$ \\
      \hline
      $\TF^2$ rejects & $p_1(1-p_2)$      & $(1-p_1)(1-p_2)$ \\
      \hline
    \end{tabular}
    }
    \captionsetup{labelformat=empty}
    \caption{Case 1: $\alpha \in \gL_1 \land \alpha \in \gL_2$}
  \end{subfigure}
  \hfill
  \begin{subfigure}[t]{0.48\linewidth}
    \centering
    \resizebox{\linewidth}{!}{%
    \begin{tabular}{|r|c|c|}
      \hline
                      & $\TF^1$ accepts   & $\TF^1$ rejects \\
      \hline
      $\TF^2$ accepts & $p_1(1-p_2)$      & $(1-p_1)(1-p_2)$ \\
      \hline
      $\TF^2$ rejects & $p_1p_2$          & $(1-p_1)p_2$ \\
      \hline
    \end{tabular}
    }
    \captionsetup{labelformat=empty}
    \caption{Case 2: $\alpha \in \gL_1 \land \alpha \notin \gL_2$}
  \end{subfigure}

  \vspace{1em}

  % Fila 2
  \begin{subfigure}[t]{0.48\linewidth}
    \centering
    \resizebox{\linewidth}{!}{%
    \begin{tabular}{|r|c|c|}
      \hline
                      & $\TF^1$ accepts       & $\TF^1$ rejects \\
      \hline
      $\TF^2$ accepts & $(1-p_1)p_2$          & $p_1p_2$ \\
      \hline
      $\TF^2$ rejects & $(1-p_1)(1-p_2)$      & $p_1(1-p_2)$ \\
      \hline
    \end{tabular}
    }
    \captionsetup{labelformat=empty}
    \caption{Case 3: $\alpha \notin \gL_1 \land \alpha \in \gL_2$}
  \end{subfigure}
  \hfill
  \begin{subfigure}[t]{0.48\linewidth}
    \centering
    \resizebox{\linewidth}{!}{%
    \begin{tabular}{|r|c|c|}
      \hline
                      & $\TF^1$ accepts       & $\TF^1$ rejects \\
      \hline
      $\TF^2$ accepts & $(1-p_1)(1-p_2)$      & $p_1(1-p_2)$ \\
      \hline
      $\TF^2$ rejects & $(1-p_1)p_2$          & $p_1p_2$ \\
      \hline
    \end{tabular}
    }
    \captionsetup{labelformat=empty}
    \caption{Case 4: $\alpha \notin \gL_1 \land \alpha \notin \gL_2$}
  \end{subfigure}

  \caption{Probabilities of all possible behaviors of $\TF^1$ and $\TF^2$}
  \label{fig:cases}
\end{figure}

$\TF$ accepts whenever $\TF^1$ or $\TF^2$ accepts.
Let us analyze case by case the probability of $\TF$ answering correctly.
%
\begin{itemize}
    \item Case 1: $p_1 + p_2 - p_1p_2 = p_1p_2 + (1-p_1)p_2 + p_1(1-p_2)$
    \item Case 2: $1 - p_2 + p_1p_2  = p_1(1-p_2) + (1-p_1)(1-p_2) + p_1p_2$
    \item Case 3: $1 - p_1 + p_1p_2  = (1-p_1)p_2 + p_1p_2 + (1-p_1)(1-p_2)$
    \item Case 4: $p_1p_2$
\end{itemize}
%
Then, if $p_1 = p_2 = 4/5$, the probability of $\TF$ answering correctly in each case is larger than 1/2.
%
\begin{itemize}
    \item Case 1: $24/25 > 1/2$
    \item Case 2: $21/25 > 1/2$
    \item Case 3: $21/25 > 1/2$
    \item Case 4: $16/25 > 1/2$
\end{itemize}

Then, if $\{\TF^1_n\}_{n \in \mathbb{N}}$ and $\{\TF^2_n\}_{n \in \mathbb{N}}$ answer correctly with probability $4/5$, the probability of $\{\TF_n\}_{n \in \mathbb{N}}$ answering correctly $\gL_1 \cup \gL_2$ is at least $1/2$.
Therefore, by the construction in Section \ref{error-reduction}, there is a family of transformers in $\RCoTtndn$ that recognizes $\gL_1 \cup \gL_2$ in $\RCoTtndn$.